\documentclass[a4paper]{article}
\usepackage[pdftex]{graphicx}
\usepackage[utf8]{inputenc}
\usepackage{enumerate}
\usepackage{icomma}
\usepackage{siunitx}
\sisetup{locale=DE}
\usepackage{tikz}
\usepackage{href-ul}
\hypersetup{
	colorlinks=true,
	linkcolor=blue,
	urlcolor=blue}
\usepackage{amssymb}
\usepackage{geometry}
\geometry{a4paper, top=15mm, left=15mm, right=15mm, bottom=15mm,
	headsep=10mm, footskip=12mm}

\begin{document}
	\begin{tikzpicture}(10,3)
		\draw[ultra thick](2,0) --(10,0) -- (10,1.3) --(2,1.3) -- (2,0);
		\draw[fill=black](2,0)-- (0,.6) -- (2,.6) -- (2,0);
		\node at (6,.6) {\large Solution sheet from \href{http://www.okuyakl.com}{www.okuyakl.com}};
	\end{tikzpicture}
	\vspace{0.5 cm}
	
	\noindent
	\begin{minipage}{0.5\textwidth}
		\noindent{\bf Task 1.}\\
		The function term is
		$$y =5^x$$
		
		\noindent{\bf Task 2.}\\
		$$ y = 100 \cdot 0.8^x$$
		
		
	\end{minipage}
	\hspace{0.5 cm}
	\begin{minipage}{0.5\textwidth}
		\noindent{\bf Task 3. a)}\\
		\includegraphics[width=7 cm]{exdrei029}
	\end{minipage}
	
	\noindent{\bf Task 3. a)}\\
	$$
	\renewcommand{\arraystretch}{2}
	\begin{array}{c|cccccccccc}
		x & 1 & 2 & 3 & 4 & 5 & 6 & 7 & 8 & 9 & 10 \\
		\hline
		y & 1.58 & 2.5 & 3.95 & 6.25 & 9.88 & 15.63 & 24.71 & 39.06 & 61.76 & 97.66 \\
	\end{array}
	$$
	
	\noindent{\bf Task 3 b.)}\\
	We also calculate the growth factor
	$$a = 2.5^{0.5} = 1.58 $$
	The increase is therefore 58%.
	
	\noindent{\bf Task 3. c)}\\
	We read: there are 50 sick people after 8.5 days.
	
	\fbox{
		\noindent
		\begin{minipage}{0.5\textwidth}
			\noindent{\bf Task 4. b)}\\
			We also calculate the growth factor
			$$a = 3^{0.5} = 1.73 $$
			The increase is therefore 73%.
			
			\noindent{\bf Task 4. c)}\\
			We read: There are 100 bacteria after 4.2 hours.
		\end{minipage}
		\hfill
		\begin{minipage}{0.4\textwidth}
			\noindent{\bf Task 4. a)}\\
			\includegraphics[width=7 cm]{exvier029}
	\end{minipage}}
	
	\noindent{\bf Task 5. a)}\\
	The function term is:
	$$y = 0.5 \cdot 1.35^x$$
	
	\noindent{\bf Task 5. b)}\\
	$$
	\renewcommand{\arraystretch}{2}
	\begin{array}{c|cccccccccc}
		x & 0 & 1 & 2 & 3 & 4 & 5 & 6 & 7 & 8 \\
		\hline
		y & 0.50 & 0.67 & 0.91 & 1.23 & 1.66 & 2.24 & 3.03 & 4.09 & 5.51
	\end{array}
	$$
	\noindent
	\begin{minipage}{0.5\textwidth}
		\noindent{\bf Task 5. b)}\\
		\includegraphics[width=7 cm]{exfunf029}
	\end{minipage}
	\hfill
	\begin{minipage}{0.5\textwidth}
		\noindent{\bf Task 5. c)}\\
		A fifth of the water surface is $\SI{4}{\meter^2}$. These will be reached after 7 days, i.e. on June 17th.
	\end{minipage}
	
	\noindent
	\begin{minipage}{0.5\textwidth}
		\noindent{\bf Task 6. a)}\\
		\includegraphics[width=7 cm]{exse029}
	\end{minipage}
	\hfill
	\begin{minipage}{0.5\textwidth}
		\noindent{\bf Task 6. b)}\\
		We also calculate the growth factor
		$$a = 0.5^{0.5} = 0.71 $$
		The decrease is therefore $100\%-71\%=29\%$.
		
		\noindent{\bf Task 6. c)}\\
		We read: after 1.5 days there are still $\SI{3}{\milli\gram}$ present.
	\end{minipage}
	\vspace{0.5 cm}
	
	\noindent{\bf Task 7.}\\
	From 2016 to 2028 there are 12 years. The growth factor is
	$$a=(1-{p \over 100})= 1-0.0125 = 0.9875$$
	The pollutant emissions after 12 years are:
	$$y=100\% \cdot 0.9875^{12} =85.99\% $$
	So it has decreased by a little over 14%. The goal is achieved.
	
	\begin{center}
		\includegraphics[width=7 cm]{../../viecher/eendcomic.pdf}
		
		Here you can go back to the \href{https://www.okuyakl.de/math/m10expL029/ae029.pdf}{task sheet}
	\end{center}
	
\end{document}